\section{Confidential-Predicate Statements and Verification}
\label{s:zkstatement}
This section gives the statement formats and the verification sequence behind Section~\ref{sec:zk}.

\emph{Policy registration.} The policy authority signs a registration (schema \code{policy\_registration\_v2}) containing the predicate identifier, the version $v$, the commitment $C_v$ of~(\ref{eq:cv}), the program identifier $I_R$ (\code{predicate\_program\_id}), the sorted list of applicable assets, the validity period, the maximum measurement age (\code{max\_measurement\_age\_s}), an optional reference to an independent reviewer's approval, and a timestamp. An independent policy reviewer may approve the parameters and the program where the assessment arrangement requires it. A commitment chosen unilaterally by a prover is not evidence that the committed policy was approved, so the verifier accepts only a $C_v$ covered by a registration that a trusted policy authority signed.

\emph{Source attestation.} The context $\kappa_i$ contains the transaction identifier, the asset identifier, the state version, and the measurement period. The source observer measures $x_i$ itself, draws $r_i$, computes $X_i$ by~(\ref{eq:xi}), and signs a statement (schema \code{source\_commitment\_attestation\_v2}) with the transaction identifier, $\kappa_i$, $X_i$, a declaration of the measurement semantics, its identity, the capture time, and its freshness limit. The source authenticates the measurement semantics; it does not act as a signing service for opaque commitments supplied by a prover, and a commitment that the source did not attest is rejected.

\emph{Public statement.} The statement consists of $C_v$, $X_i$, $\kappa_i$, $D_i$, $v$, the predicate identifier, and $q_i$. $D_i$ is SHA-384 over the action domain label and the canonical encoding of the proposal; it is the same digest that the authorization records as its target configuration digest and that the receipt records as its proposal digest.

\emph{Verification sequence.} The verifier of the bound statement performs the following checks in order and rejects at the first failure.
\begin{enumerate}
\item The registration carries a valid signature of a trusted policy authority at its timestamp; its commitment, predicate identifier, and version equal those of the statement; it is valid at verification time; it covers the asset of $\kappa_i$; and it sets a non-negative integer maximum measurement age.
\item The source attestation carries a valid signature of a trusted source observer at its capture time; its commitment equals $X_i$ and its context equals $\kappa_i$; and the elapsed time since capture is non-negative and at most the smaller of the attested freshness limit and the registered maximum age.
\item The authorization carries a valid signature of a trusted authorizer whose key is valid at the authorization's timestamp, has the schema \code{authorization\_v2} and the decision ``authorized,'' and names the transaction and asset of $\kappa_i$ and the target configuration digest $D_i$.
\item If a proof is supplied, the registration must name $I_R$; the receipt must be succinct, lifted from a single segment of at most $2^{\ZkMaxPo}$ cycles, and valid for $I_R$; and its public output must equal the verifier's canonical encoding of the statement byte for byte. Otherwise, if an opening is supplied, the relation must hold with the witness in the clear (reviewer path). If neither is supplied, the statement is rejected.
\end{enumerate}
Every public field used to support a claim is thus either constrained by the relation or checked through a bound authenticated statement. The authorization must match the transaction, asset and action digest of the statement, but not the state version in $\kappa_i$, and the archive verifier checks that the receipt binds the same action digest (obligation~(ii), Section~\ref{sec:semantic}).

\emph{Program binding.} A proof for general computation establishes only that the program against which it is verified produced a given output. If the verifier accepted a program named by the prover, the prover could present a valid proof for a program that emits a well-formed statement with $q_i=1$ irrespective of $\theta$, or that uses a different predicate, encoding, or commitment function. Placing $I_R$ in the signed registration makes the choice of program, and with it the evaluated $Q$, the commitment function, and the canonical encoding, a decision of the policy authority rather than of the prover. Within the program, $Q$ is selected from a fixed catalog by the predicate identifier, which the registration also covers, so the pair of identifiers determines the evaluated predicate. The binding fixes the form of $Q$; its parameters remain hidden behind $C_v$. $I_R$ is the image identifier of the compiled guest program. Reproducing it from source requires the same guest toolchain, so a reviewer who approves a registration either rebuilds the program with that toolchain or relies on the party that built it (Section~\ref{sec:limits}).

\emph{Predicate scope.} Relation~(\ref{eq:rel}) defines a Boolean policy check and permits multiple acceptable actions. It verifies the bound predicate; it does not reproduce the reasoning of the AI advisor. The catalog contains one predicate, a rotation threshold: $q_i$ holds if the measured key age in days reaches a confidential minimum and the target profile of the action belongs to a confidential permitted set. Temporal predicates would require authenticated prior state and a corresponding extension of the relation.

\emph{Proving environment.} The host constructs the vendor's local prover explicitly, so the environment variable that selects a prover has no effect; the Python backend additionally removes the variables that enable the development mode, select another prover, or name an external prover binary from the host's environment, and the host refuses to run when the development-mode variable is set. The development mode is also disabled at compile time. The witness reaches the host through a file in a temporary directory that is deleted after proving.

\section{Commitment Construction}
\label{s:zkcommit}
The reference implementation and the guest program instantiate the commitment $\Com$ of~(\ref{eq:cv}) and~(\ref{eq:xi}) as
\begin{multline*}
\Com(m; r) = \mathrm{SHA\text{-}384}(\texttt{"ABD/v2/commitment"}\\
\|\, \texttt{0x00} \,\|\, \mathrm{u16be}(|r|) \,\|\, r \,\|\, m),
\end{multline*}
where $r$ consists of 32 bytes from the operating system's random number generator, drawn afresh for every commitment, $\mathrm{u16be}$ is a two-byte big-endian length, and $m$ is the canonical encoding of the object with fields $v$ and $\theta$ for~(\ref{eq:cv}) or $\kappa_i$ and $x_i$ for~(\ref{eq:xi}). The domain label and separator are fixed, $r$ has a fixed, length-prefixed size, and $m$ is the remainder, so the input encoding is unambiguous and binding reduces to the collision resistance of SHA-384. Collision resistance alone is not a general proof of hiding. Hiding holds when SHA-384 is modeled as a random oracle; it is a stated assumption, not a proven property of the hash function.

\section{Cross-Language Encoding Check}
\label{s:zkencoding}
This section supports the byte-for-byte comparison of the program output with the verifier's encoding that Section~\ref{sec:zk} requires. The verifier compares the public output of the receipt with its own Python encoding of the statement, while the guest program uses an independently written Rust encoder. Both accept the same value space (printable-ASCII strings, which exclude line terminators; Booleans; null; integers within $\pm(2^{53}-1)$; arrays; and objects with printable-ASCII string keys, which both emit in sorted order), and both reject floating-point values. Equality of the two encoders was checked by executing the guest without proving on vectors containing quotes and backslashes, negative integers, null, empty containers, and nested objects, and on the proven statement; every output equaled the Python encoding. The Python evaluator and the guest apply the same strict types to the predicate inputs: a Boolean is not accepted as an integer, a non-string entry of the permitted set never matches, and an off-schema input makes the Python evaluator raise an error and the guest abort, so that no proof exists for it. Regression tests cover these cases.

\section{Disclosure Analysis of the Confidential Predicate}
\label{s:zkprivacy}
This section expands the disclosure remark at the end of Section~\ref{sec:zklimits}. Zero knowledge, where it holds, limits additional witness disclosure beyond the public statement and other available observations. The statement itself discloses $q_i$, $\kappa_i$, $v$, the predicate identifier, and $D_i$, and $D_i$ identifies an action that the audit record discloses. For the rotation-threshold predicate, a statement with $q_i=1$ reveals that the target profile is in the permitted set and that the minimum does not exceed the hidden key age, which bounds the minimum wherever the key age is independently known; a statement with $q_i=0$ reveals that at least one condition failed. A sequence of results for varied actions or periods can therefore narrow both parameters. The construction hides parameters, not the form of $Q$: the structure of the predicate is known to anyone who can inspect the program that $I_R$ identifies.

As Section~\ref{sec:zklimits} states, the control identifier of a succinct receipt discloses a coarse execution length. This holds although the vendor's security model states that recursion proofs leak no information about execution length~\cite{risc0sec}: the measured control identifier for the evaluated statement names the lift program for $2^{\ZkPo}$ cycles. The cycle count can depend on $\theta$, for example through the length of the permitted set that the guest scans. The single-segment bound of the verifier limits accepted receipts to lift programs of at most $2^{\ZkMaxPo}$ cycles. A composite receipt carries one proof per segment~\cite{risc0}; succinct receipts are selected because they are hash-based and consist of one recursion proof regardless of the segment count, not for hiding the length.

The receipt bytes contain neither the canonical encodings of $\theta$ and $x_i$ nor $r$ and $r_i$, in raw or hexadecimal form. This check excludes verbatim inclusion of the witness; it is not evidence of zero knowledge.

The result, action, timing, and released rationale can reveal information about policy. Disclosure controls therefore use audience-specific allowlists, reviewed rationale templates, and explicit leakage analysis; free-text fields are not structurally incapable of carrying secrets. Secret copies, backups, operator access, and any reviewer escrow belong to the confidentiality boundary; a commitment does not remove these copies.

Fresh commitments and secure removal of obsolete witnesses reduce the exposure of stored openings. They do not provide semantic forward secrecy for policy parameters that remain unchanged across epochs, and no bound expressed as a fraction of historical policy content follows from re-randomization alone.

\section{Proof-System Assumptions}
\label{s:zkpq}
This section states the assumptions behind the soundness levels and the zero-knowledge caveat of Section~\ref{sec:zklimits}. A post-quantum assurance profile requires a proof system whose construction and parameters support the intended quantum adversary model. STARKs provide a relevant foundation~\cite{stark}. Chiesa, Manohar, and Spooner prove that succinct non-interactive arguments obtained from interactive oracle proofs with round-by-round soundness are unconditionally sound in the quantum random oracle model, and obtain from Micali's construction a zero-knowledge succinct argument secure in that model~\cite{cms}. These results concern idealized constructions; they do not establish the security of a particular implementation, its parameters, its hash functions, or its zero-knowledge variant, and the name of a proof-system family does not establish a concrete security level.

RISC Zero's security model states conjectured levels of 99 bits for the recursion prover and 96 bits for the RISC-V prover, describes both as quantum safe, and bases both on the random oracle model and on the conjecture that the best known attack on its STARK is the best possible attack~\cite{risc0sec}. Neither figure follows from a security reduction, and neither is an end-to-end level for a receipt. The levels in Section~\ref{sec:zklimits} come from the vendor's own soundness calculator (\code{risc0\_zkp::prove::soundness} in risc0-zkp 3.0.5, the version pinned by version \ZkVersion\ of the RISC Zero zkVM~\cite{risc0}), which the host evaluates for every segment size from $2^{14}$ to $2^{22}$ cycles and for the recursion circuit. The toy-model level assumes the ethSTARK toy-problem conjecture, the vendor's default; the strict level assumes a list-decoding conjecture for the proximity test; the proven level assumes neither. A regression test checks the calculator against the vendor's published reference value and checks that the level does not increase with the segment size. A succinct receipt is sound only if the segment proof that it verifies is sound, so its end-to-end level is the minimum of the segment and recursion levels, and the bound on the accepted segment size fixes the worst case: the verifier derives the segment size from the control identifier and rejects join programs and larger lift programs. Because the level falls as the segment size grows, the worst accepted segment (\ZkSegBits~bits) lies above the vendor's 96-bit figure for the RISC-V prover, to which the calculator falls only for larger segments.

Wrapping a STARK receipt in a Groth16 proof yields a smaller receipt, but its security rests on pairings over the BN254 curve, a knowledge-of-exponent assumption, and a trusted setup, and it is not post-quantum~\cite{risc0sec}. Such receipts are neither produced nor accepted.

Zero knowledge for STARKs is subtle. Hab\"ock and Al~Kindi analyze zero knowledge for univariate argument systems that use the fast Reed--Solomon interactive oracle proof of proximity (FRI) as polynomial commitment~\cite{habock}. Citing that analysis, a 2024 vendor advisory for zkVM versions up to 1.0.2 stated that several STARK implementations, including the zkVM, did not meet the requirements for asserting zero knowledge provably~\cite{risc0adv}. The current security model states that the zkVM targets perfect zero knowledge but that no mathematical argument for it has been written~\cite{risc0sec}. Zero knowledge is therefore a stated assumption of this work; soundness does not depend on it.
